\documentclass[10pt,a4paper,serif,expansion=false]{moderncv}



% ---------- Theme ----------
\moderncvstyle{classic}   % casual | classic | banking | oldstyle | fancy
\moderncvcolor{blue}      % blue | orange | green | red | purple | grey | black
\usepackage[scale=0.75]{geometry}
\usepackage{hyperref}
\hypersetup{
    colorlinks=true,
    linkcolor=blue,
    urlcolor=blue,
    citecolor=blue
}

% Width of the left ("date") column
\setlength{\hintscolumnwidth}{2.6cm}
\AtBeginDocument{\recomputecvlengths}

% ---------- Personal information ----------
\name{}{Shin JeongBean}
\title{Curriculum Vitae}
\email{shinjeongbean@gmail.com}
\extrainfo{B10104030@ntu.edu.tw}
\homepage{citron2357.github.io}          % personal website, if any
\extrainfo{ORCID: 0009-0004-7856-142X \quad \href{https://philpeople.org/profiles/jeongbean-shin}{PhilPeople profile}}  % ORCID, if any

\begin{document}
\makecvtitle

% ============================================================
\section{Education}
\cventry{2021\textbf{--}present}{B.A.\ in Philosophy}{National Taiwan University}{Taipei, Taiwan}{\textit{GPA: 4.14/4.3} (3.97/4.0, with A+ counted as 4.0)}{Expected graduation: June 2027. Leave of absence for mandatory military service in South Korea, Oct 2023 \textbf{--} Apr 2025.}

% ============================================================
\section{Areas of Interest}
\cvitem{}{Metaphysics, Metametaphysics, Philosophical Logic, Philosophy of Art}

% ============================================================
\section{Presentations}
\cvitem{Nov. 2026}{(forthcoming) \emph{Belief Revision in Truthmaker Semantics} (refereed). Taiwan Philosophical Association Annual Conference 2026, National Chung Cheng University, Chiayi, Taiwan.}
\cvitem{Apr. 2026}{\emph{Essentialism and the Essence of what it is to be Necessary} (refereed). TAPSC 2026, Soochow University, Taipei, Taiwan. (A version is under review at \emph{Philosophical Studies}.)}

% ===========================================================
\section{Commentary}
\cvitem{Apr. 2026}{Commentator on Lin Siao-Ting, ``A Fragmentalist Passage of Time''. TAPSC 2026, Soochow University, Taipei, Taiwan.}

% ============================================================
\section{Teaching Experience}
\cventry{Fall 2026}{Teaching Assistant \textmd{(in progress)}}{National Taiwan University}{PHL1921 \emph{Using AI to Tackle Philosophical Questions about Truth in Fiction}}{Instructor: Prof.\ Hsien-Chieh Nihel Jhou.}
  {\footnotesize \textcolor{color1}{General-education course, taught in English. Deliver a short lecture and lead discussion in most of the weekly three-hour sessions; grading group presentations and the final exam.}}
% ============================================================
\section{Works in Progress}
% TODO: add advisor name to the thesis entry if desired.
\cventry{}{\textmd{\emph{Essentialism and the Essence of what it is to be Necessary}}}{}{Under review at \emph{Philosophical Studies}, Springer. \href{https://github.com/citron2357/Ess-Nec/blob/main/main.pdf}{[manuscript]}}{}
  {\footnotesize \textcolor{color1}{I argue that if necessity is grounded in essence, the necessity of essence should obtain essentially of what it is to be necessary.}}
\cventry{}{\textmd{\emph{Belief Revision in Truthmaker Semantics} \href{https://github.com/citron2357/Logic-of-belief-revision-with-truthmaker-semantics/blob/main/main.pdf}{[manuscript]}}}{}{}{}
  {\footnotesize \textcolor{color1}{I develop several hyperintensional AGM-style belief revision frameworks with the resources from Finean truthmaker semantics to model local rationality.}}
\cventry{}{\textmd{\emph{Demystifying Anchoring}}}{}{B.A.\ thesis, National Taiwan University. Advisor: Prof.\ Duen-Min Deng. \href{https://github.com/citron2357/Demystifying-Anchoring/blob/main/main.pdf}{[unfinished manuscript]}}{}
  {\footnotesize \textcolor{color1}{I argue that anchoring need not be a \textit{sui generis} metaphysical relation by examining various possible cases of exportation.}}
\cventry{}{\textmd{\emph{Abstract Object Theory and the Fictionally Fictional}}}{}{NSTC Research Project for College Students. Advisor: Prof.\ Hsien-Chieh Nihel Jhou. \href{https://github.com/citron2357/AOTandFF/blob/main/Main.pdf}{[unfinished manuscript]}}{}
  {\footnotesize \textcolor{color1}{I present a counterargument against Zalta's abstract object theory from the case of fictional fictions and their objects.}}

% ============================================================
\section{Graduate Coursework}
\cvitem{Fall 2026}{PHL~7954, Hyperintensionalism (in progress)}
\cvitem{Spring 2026}{PHL~7763, Epistemic Logic; PHL~7599, The Philosophy of Bertrand Russell, 1896--1900; PHL~7952, Zeno's Paradoxes; PHL~7558, Wittgenstein's \emph{Philosophical Investigations}, \S\S\,185--693}
\cvitem{Fall 2025}{PHL~7598, Essence and Essentialism; PHL~7557, Wittgenstein's \emph{Philosophical Investigations}, \S\S\,1--184; PHL~7950, Machine Ethics}

% ============================================================
\section{Scholarships and Awards}
\cvitem{2026\textbf{--}27}{National Science and Technology Council (NSTC) Research Project for College Students (research grant), Jul 2026 \textbf{--} Feb 2027.}
\cventry{2026}{\textmd{National Taiwan University System President's Award}}{}{}{}{\footnotesize \textcolor{color1}{Newly established in 2026; awarded to international students at NTU, NTNU, and NTUST who received the Academic Achievement Award at least twice within four years.}}
\cvitem{2023}{The 13th National Taiwan University Student Laureate for Philosophical Treatise, 4th prize.}
\cvitem{2022, 2026}{NTU Scholarship for International Degree Students}
\cventry{2021\textbf{--}23, 2025\textbf{--}26}{\textmd{Academic Achievement Award}}{}{}{{\small 2021 Fall; 2022 Fall; 2023 Spring; 2025 Fall; 2026 Spring}}{\footnotesize \textcolor{color1}{Awarded each semester to approximately the top 5\% of students in the department.}}

% ============================================================
\section{Academic Activities}
\cventry{2026\textbf{--}present}{NTU Logic}{Founder and Organizer}{}{}{\footnotesize \textcolor{color1}{Weekly reading group founded with 9 participants from philosophy and mathematics (undergraduate and graduate). Texts: Boolos et al.\ (2007), \emph{Computability and Logic}; Priest (2008), \emph{An Introduction to Non-Classical Logic}; Bacon (2023), \emph{A Philosophical Introduction to Higher-Order Logics}}}
\cventry{2025\textbf{--}present}{Philosophy of Music Study Group}{Founder and Organizer}{}{}{\footnotesize \textcolor{color1}{Interdisciplinary reading group of 6 regular members (backgrounds in philosophy, art history, piano performance, and computer science), meeting biweekly on average; each session is devoted to a selected paper/book chapter in the philosophy of music.}}
\cventry{Spring 2026}{Category Theory Study Group}{Participant}{}{}{{\footnotesize \textcolor{color1}{Biweekly reading group of undergraduate and graduate students; text: Smith (forthcoming), \emph{Category Theory: A Gentle Introduction.}}}}


% ============================================================
\section{Skills}
\cvitem{Languages}{Korean (native); English (fluent: TOEFL iBT: 5.5/6  (0\textbf{--}120 scale: 110), July 2026); Mandarin Chinese (fluent: TOCFL C2, 2020); Classical Chinese (reading, upper-intermediate)}
\cvitem{Computer}{LaTeX; Prolog}
\cvitem{Music}{Training in piano and music composition (2015\textbf{--}2020); coursework in chamber music at National Taiwan Normal University}

\end{document}